\chapter{استدلال و برهان}
\thispagestyle{empty}

\begin{quote}
	{\lotus
		سعی کنید آنچه را شهودی به نظر می‌رسد، به‌طور رسمی و دقیق اثبات کنید و آنچه را که به‌طور رسمی و دقیق اثبات کرده‌اید، به‌طور شهودی درک کنید. این یک ورزش مغزی جالب است. \quad - پولیا%
		\LTRfootnote{ George Polya}
	}
\end{quote}

یک برهان در ریاضیات اغلب شامل مجموعه‌ای منطقی از گام‌هاست که درستی یک گزاره‌ی عمومی را تأیید می‌کند.

\section{استنتاج}
استدلال استنتاجی روش نتیجه‌گیری با استفاده از حقایقی است که درستی آن‌ها را پذیرفته‌ایم. وقتی از استدلال استنتاجی استفاده می‌کنیم، مطمئن هستیم که نتیجه همیشه درست است. قضایای کلی احکامی هستند که همیشه برقرار می‌باشند. به این نوع استدلال، اثبات مستقیم هم می‌گویند.

\begin{example}
	ثابت کنید که مجموع یک عدد صحیح مثبت فرد و یک عدد صحیح مثبت زوج همیشه عددی فرد است. \\
	\textbf{پاسخ:} \\
	یک عدد صحیح مثبت فرد به صورت
	$2k-1$
	نوشته می‌شود که
	$k \in \mathbb{Z}^+$. \\
	یک عدد صحیح مثبت زوج به صورت
	$2k$
	نوشته می‌شود که
	$k \in \mathbb{Z}^+$. \\
	فرض کنید $m$ و $n$ به ترتیب دو عدد صحیح مثبت فرد و زوج باشند.
	$$m=2p-1 \ , \quad n=2s \ , \quad p,s \in \mathbb{Z}^+$$
	$$\Rightarrow m+n=2p-1+2s=2(p+s)+1$$
	چون
	$p+s \in \mathbb{Z}^+$،
	پس طبق عبارت بالا
	$m+n$
	فرد است.
	\qed
\end{example}

\begin{example}
	ثابت کنید که اگر مجموع ارقام عددی چهاررقمی بر ۳ بخش‌پذیر باشد، آنگاه خود آن عدد نیز بر ۳ بخش‌پذیر است. \\
	\textbf{پاسخ:}
	فرض کنید $n$ عددی چهاررقمی به این صورت باشد:
	$n = a_3 a_2 a_1 a_0$.
	داریم:
	\[
	\begin{cases}
		a_3 = p \times 10^3 \\
		a_2 = q \times 10^2 \\
		a_1 = r \times 10^1 \\
		a_0 = s \times 10^0
	\end{cases}
	\quad
	0 \leq p, q, r, s \leq 9 \quad \text{و} \quad p \neq 0
	\]
	$$p+q+r+s = 3k \ , \ k \in \mathbb{Z}$$
	$$\Rightarrow n = p \times 10^3 +  q \times 10^2 + a_1 + r \times 10^1 + s \times 10^0$$
	$$n = p \times 10^3 +  q \times 10^2 + a_1 + r \times 10^1 + (3k-p-q-r)$$
	$$n = (p \times 10^3 -p) + (q \times 10^2 -q) + (r \times 10^1 -r) +3k$$
	$$n = p(10^3-1) + q(10^2-1) + r(10^1-1) +3k$$
	$$n = 999p +99q +9r +3k = 3(333p+33q+3r+k)$$
	از آن‌جا که
	$333p+33q+3r+k \in \mathbb{Z}$،
	نتیجه می‌شود که $n$ بر ۳ بخش‌پذیر است.
	\qed
\end{example}

\section{استقرای ریاضی}

\begin{theorem}
	فرض کنید
	$P(n)$
	حکمی درباره‌ی عدد طبیعی $n$ باشد. اگر
	$P(1)$
	درست باشد و از درستی
	$P(k)$
	درستی
	$P(k+1)$
	نتیجه شود، در این صورت
	$P(n)$
	برای هر عدد طبیعی $n$ نیز درست است.
\end{theorem}

برای استفاده از استقرای ریاضی، گام‌های زیر را برمی‌داریم: \\
گام ۱: درستی حکم را برای
$n=1$
نشان می‌دهیم. \\
گام ۲: ثابت می‌کنیم که اگر حکم برای
$n=k$
درست باشد، آنگاه برای
$n=k+1$
نیز درست است. \\
آنگاه نتیجه می‌شود که حکم برای هر عدد طبیعی $n$ درست است.

\begin{theorem}[استقرای تعمیم‌یافته]
	فرض کنید
	$P(n)$
	حکمی درباره‌ی عدد طبیعی $n$ باشد. اگر
	$P(m)$
	برای
	$m>1$
	درست باشد و از درستی
	$P(k)$
	برای هر عدد طبیعی
	$k \geq m$
	درستی
	$P(k+1)$
	نتیجه شود، آنگاه
	$P(n)$
	برای هر عدد طبیعی
	$n \geq m$
	درست است.
\end{theorem}

برای استفاده از استقرای تعمیم‌یافته، گام‌های زیر را برمی‌داریم: \\
گام ۱: $m$ مناسب را به دست می‌آوریم. \\
گام ۲: درستی حکم را برای $n=m$ نشان می‌دهیم. \\
گام ۳: ثابت می‌کنیم که اگر حکم برای
$n=k \geq m$
درست باشد، آنگاه حکم برای
$n=k+1$
نیز درست است. \\
آنگاه نتیجه می‌شود که حکم برای هر عدد طبیعی
$n \geq m$
درست است.

\section{مثال نقض}
به مثالی که نشان دهد نتیجه‌گیری کلی غلط است، مثال نقض می‌گویند.
\begin{example}
	نشان دهید که گزاره‌های زیر درست نیستند. \\
	۱. اگر
	$n \in \mathbb{Z}$
	و
	$n^2$
	بر ۴ بخش‌پذیر باشد، آنگاه $n$ بر ۴ بخش‌پذیر است. \\
	۲. 
	\\ \textbf{پاسخ:} \\
	۱. عدد
	$n=6$
	را در نظر بگیرید.
	$n^2=36$
	بر ۴ بخش‌پذیر است ولی
	$n=6$
	بر ۴ بخش‌پذیر نیست. \\
	۲.
\end{example}

\section{اثبات برگشت‌پذیر}
گاهی برای اثبات بعضی از قضیه‌ها - به‌خصوص درمورد تساوی‌ها، نامساوی‌ها و تساوی مجموعه‌ها - با استفاده از درستی یک حکم به یک رابطه‌ی بدیهی و یا فرض قضیه می‌رسیم. در چنین حالتی، برای تکمیل اثبات باید نشان دهیم که تمام مراحل انجام‌شده بازگشت‌پذیر هستند وگرنه درستی اثبات تأیید نمی‌شود.

\section{عکس نقیض و برهان خلف}
اثبات به کمک عکس نقیض و اثبات به کمک برهان خلف بسیار شبیه یکدیگرند. \\
در اثبات به کمک عکس نقیض از معادل بودن دو گزاره‌ی زیر استفاده می‌کنیم:
$$(p \rightarrow q) \Leftrightarrow ((\neg q) \rightarrow (\neg p))$$

\begin{example} \label{prexn2even}
	فرض کنید $n$ عددی طبیعی باشد. ثابت کنید که اگر
	$n^2$
	زوج باشد، آنگاه $n$ زوج است. \\
	\textbf{پاسخ:}
	کافی است ثابت کنیم که اگر $n$ زوج نباشد (فرد باشد)، آنگاه
	$n^2$
	زوج نیست (فرد است). با فرض این‌که $k$ عدد صحیح نامنفی باشد، داریم:
	$$n=2k+1 \Rightarrow n^2=(2k+1)^2=4k^2+4k+1=2(2k^2+2k)+1=2k'+1$$
	از آن‌جا که $k'$ عددی صحیح و نامنفی است،
	$n^2$
	فرد است.
	بنابراین اگر $n$ فرد باشد،
	$n^2$
	فرد است. در نتیجه اگر
	$n^2$
	زوج باشد، آنگاه $n$ زوج است.
	\qed
\end{example}
\bigskip
برهان خلف به زبان ریاضی به صورت زیر است:
$$(p \rightarrow q) \Leftrightarrow ((p \wedge (\neg q)) \rightarrow \bot)$$
برای استفاده از برهان خلف گام‌های زیر را در نظر می‌گیریم: \\
گام ۱: فرض می‌کنیم نتیجه‌ی مطلوب درست نباشد. \\
گام ۲: نشان می‌دهیم که این فرض نتیجه‌ای به دست می‌دهد که حقایق دانسته‌شده را نقض می‌کند. \\
گام ۳: حال که به یک تناقض رسیده‌ایم، معلوم می‌شود که فرضی که در گام اول کرده بودیم نادرست است. بنابراین نتیجه‌ی مطلوب باید درست باشد.

\begin{example}
	نشان دهید که
	$\sqrt{2}$
	عددی گنگ است. \\
	\textbf{پاسخ:}
	فرض کنید که
	$\sqrt{2}$
	گنگ نیست، یعنی گویا است. در نتیجه فرض کنید
	$\sqrt{2}=\frac{m}{n}$
	که
	$m,n \in \mathbb{Z}$
	و همچنین $m$ و $n$ عامل مشترکی ندارند.
	$$\sqrt{2}=\frac{m}{n} \Rightarrow 2=\frac{m^2}{n^2} \Rightarrow m^2=2n^2$$
	که این یعنی $m$ عددی زوج است (بنابر مثال
	\ref{prexn2even}).
	$$m=2k \Rightarrow m^2=4k^2 \Rightarrow 2n^2=4k^2 \Rightarrow n^2=2k^2$$
	که این یعنی $n$ نیز عددی زوج است (بنابر مثال
	\ref{prexn2even}). \\
	اما فرض کرده بودیم که $m$ و $n$ هیچ عامل مشترکی ندارند و اکنون به این نتیجه رسیدیم که هر دو عامل مشترک $2$ را دارند. بنابراین، نمی‌توان
	$m,n \in \mathbb{Z}$
	پیدا کرد که عامل مشترک نداشته باشند و 
	$\sqrt{2}=\frac{m}{n}$
	عدد گویا باشد. بنابراین،
	$\sqrt{2}$
	عددی گنگ است.
	\qed
\end{example}

\begin{example}
	ثابت کنید که اگر
	$m,n \in \mathbb{Z}$،
	آنگاه
	$m^2-4n-7 \neq 0$.
	\\ \textbf{پاسخ:}
	فرض کنید که
	$m^2-4n-7=0$.
	$$\Rightarrow m^2=4n+7 \Rightarrow m^2=4n+6+1=2(2n+3)+1$$
	از آن‌جا که
	$2n+3 \in \mathbb{Z}$،
	در نتیجه
	$m^2$
	عددی فرد است که یعنی $m$ نیز عددی فرد است. (می‌توان مشابه مثال
	\ref{prexn2even}
	ثابت کرد.) فرض کنید
	$m=2p+1$
	که
	$p \in \mathbb{Z}$.
	$$(2p+1)^2-4n-7=0 \Rightarrow 4p^2+4p+1-4n-7=0 \Rightarrow 2p^2+2p-2n=3 \Rightarrow2(p^2+p-n)=3$$
	از آن‌جا که
	$p,n \in \mathbb{Z}$،
	تساوی آخر نمی‌تواند درست باشد. بنابراین اگر
	$m,n \in \mathbb{Z}$،
	آنگاه
	$m^2-4n-7 \neq 0$.
	\qed
\end{example}

\section{مسائل}

\begin{problem}
	ثابت کنید که حاصل‌ضرب دو عدد فرد، همواره عددی فرد است.
\end{problem}

\begin{problem}
	ثابت کنید که اگر
	$n \in \mathbb{Z}^+$،
	آنگاه
	$2^{n+2}+3^{2n+1}$
	بر ۷ بخش‌پذیر است.
\end{problem}

\begin{problem}
	ثابت کنید که
	$\sqrt{2} + \sqrt{3}$
	عددی گنگ است.
\end{problem}

\begin{problem}
	ثابت کنید که
	$\sqrt[5]{2}$
	عددی گنگ است.
\end{problem}

\begin{problem}
	ثابت کنید که به‌ازای هر
	$a,b \in \mathbb{Z}$،
	داریم
	$12a^2-6b^2 \neq 0$.
\end{problem}
